% 第五部分：后续工作计划
\section{后续工作计划}

\begin{frame}{已完成工作总结}
    \begin{columns}[T]
        \begin{column}{0.48\textwidth}
            \begin{block}{\normalsize {理论工作}}
                \begin{itemize}
                    \item[\textcolor{accentgreen}{\checkmark}] 完成方法整体框架设计
                    \item[\textcolor{accentgreen}{\checkmark}] 完成 MDP 问题形式化
                    \item[\textcolor{accentgreen}{\checkmark}] 完成表格型 Q-Learning 算法设计
                    \item[\textcolor{accentgreen}{\checkmark}] 完成收敛性理论分析
                \end{itemize}
            \end{block}
        \end{column}

        \begin{column}{0.48\textwidth}
            \begin{block}{\normalsize {实验工作}}
                \begin{itemize}
                    \item[\textcolor{accentgreen}{\checkmark}] 完成数据预处理流程
                    \item[\textcolor{accentgreen}{\checkmark}] 完成基线方法实现
                    \item[\textcolor{accentgreen}{\checkmark}] 完成对比实验
                    \item[\textcolor{accentgreen}{\checkmark}] 完成消融实验
                    \item[\textcolor{accentgreen}{\checkmark}] 完成统计显著性检验
                \end{itemize}
            \end{block}
        \end{column}
    \end{columns}
\end{frame}

\begin{frame}{后续进度安排}
    \begin{table}
        \centering
        \begin{tabular}{cllc}
            \toprule
            \textbf{阶段} & \textbf{时间} & \textbf{工作内容} & \textbf{预期成果} \\
            \midrule
            \rowcolor{lightblue}
            一 & 第 1-2 周 & 补充实验 & …… \\
            二 & 第 3-4 周 & 完善论文撰写 & 完成初稿 \\
            \rowcolor{lightblue}
            三 & 第 5-6 周 & 论文修改与优化 & 完成二稿、三稿 \\
            四 & 第 7-8 周 & 格式审查与定稿 & 提交终稿 \\
            \rowcolor{lightblue}
            五 & 第 9-10 周 & 答辩准备 & PPT、演示材料 \\
            \bottomrule
        \end{tabular}
    \end{table}

    \vspace{0.5em}

    \begin{exampleblock}{\normalsize 里程碑节点}
        \begin{itemize}
            \item \textbf{\color{accentgreen}第 2 周}：完成所有补充实验
            \item \textbf{\color{accentgreen}第 4 周}：完成毕业论文初稿
            \item \textbf{\color{accentgreen}第 8 周}：完成论文定稿
        \end{itemize}
    \end{exampleblock}
\end{frame}

\begin{frame}{预期成果}
    \begin{columns}[T]
        \begin{column}{0.48\textwidth}
            \begin{block}{\normalsize {学术成果}}
                \begin{itemize}
                    \item 完成毕业论文一篇
                    \item 发表学术论文 1 篇 (在准备)
                    \item 开源代码库 (计划)
                \end{itemize}
            \end{block}
        \end{column}

        \begin{column}{0.48\textwidth}
            \begin{block}{\normalsize {实践成果}}
                \begin{itemize}
                    \item 完整的 BCI 时间窗优化系统
                    \item 可复现的实验代码
                    \item 详尽的实验分析报告
                \end{itemize}
            \end{block}
        \end{column}
    \end{columns}
\end{frame}


% 结束页
\section{致谢}

\begin{frame}
    \begin{tikzpicture}[remember picture,overlay]
        \fill[top color=primaryblue,bottom color=secondaryblue]
            (current page.south west) rectangle (current page.north east);
    \end{tikzpicture}
    \vfill
    \centering
    {\Large\bfseries\color{white} 感谢各位老师的聆听与指导！}

    \vspace{2em}

    {\Large\color{white} 恳请各位老师批评指正}

    \vspace{2em}

    {\color{white}\rule{0.5\textwidth}{1pt}}

    \vspace{0.5em}

    {\small\color{white} 姓名：邓凯璇 \quad|\quad 学号：37220222203575 \quad|\quad 邮箱：3087958840@qq.com}
    \vfill
\end{frame}

\begin{frame}
    \vfill
    \centering
    {\Large\bfseries\color{primaryblue} Q \& A}

    \vspace{0.5em}

    {\Large\color{darkgray} 请各位老师提问}
    \vfill
\end{frame}
